% !PRINTED 480
\weberpage{480}
In the same way the class invariants listed in the table at the end
have been computed for
\[
m=75,\quad 36,\quad 100,\quad 63,\quad 175
\]
and these computations can also be continued still further.

II. The \gesperrt{Schlaefli} modular equations can be turned to account in
various other ways as well for the computation of class
invariants.

If we set
\[
\omega=-\frac{m}{\omega},\quad \omega=\sqrt{-m},
\]
then
\weq{11}{f(\omega)=f\Bigl(\frac{\omega}{m}\Bigr)=f(\sqrt{-m});}
if therefore in the modular equation for the transformation degree $m$
\[
u=v
\]
is set, there results an equation among whose roots
$u=f(\sqrt{-m})$ occurs. In the system I. of § 73 we must then
always set $A=2$, and the formulae result:
\[
m=3,\quad B=u^6-\frac{8}{u^6}=2,
\]
\[
m=5,\quad B=u^4-\frac{4}{u^4}=2,
\]
\[
m=7,\quad B=u^6+\frac{8}{u^6}=9,
\]
\[
m=11,\quad B=u^2-\frac{2}{u^2},\quad B^5-B^3-2B=2,
\]
\[
m=13,\quad B=u^{12}-\frac{64}{u^{12}}=9\cdot2^5,
\]
\[
m=17,\quad B=u^8+\frac{16}{u^8},\quad B^2-68B-544=0,
\]
\[
m=19,\quad B=u^6-\frac{8}{u^6},\quad B^3-38B^2+252B-648=0.
\]
The cases $m=3,7,11,19$ yield no new results, but can be used
for the verification of those found earlier;
from $m=5,13$ we obtain by the solution of a
quadratic equation:
%CONT>>
